% TODO: Add much much more!
\subsection{Motion in Space}

\content{

}{% presentation
\vspace{3in}
}{% nopresentation
}{% comments
demonstrate how the derivative of a vector valued function is the tanged vector,
which is the direction of motion of a particle traveling along $\vec{r}(t)$. 

Similarly, the second derivative is the acceleration, just as in the calc I
case. 

Mention that the speed of the particle at a point is $|\vec{r}'(t)|$.
}

\content{
\begin{example}
A moving particle starts at an initial position of $\vec{r}(0) = \la 1,0,0\ra$,
and with initial velocity $\vec{v}(0) = \la 1,-1,1\ra$.  It's acceleration is
given by $\vec{a}(t) = \la 4t, 6t, 1 \ra$. Find its velocity and position
functions. 
\end{example}
}{% presentation
\vspace{4in}
}{% nopresentation
}{% comments
Mention that there are good applications of vector-valued functions to
projectile motion and Kepler's Laws of Planetary Motion in the book if you are
interested in reading them. 
}

